\documentclass{beamer}
\usepackage[T1]{fontenc}
\usepackage[utf8]{inputenc}
\usepackage[brazilian]{babel}
\usepackage{hyperref}
\usepackage{listingsutf8}
\usepackage{color}
\usepackage{theoremref}
\usepackage{verbatim,amssymb,latexsym,amsmath,mathrsfs}

\usecolortheme{rose}
\usetheme{Hannover}

\hypersetup
{
  colorlinks=true,
  urlcolor=blue,
  citecolor=blue,
  linkcolor=blue
}

\definecolor{darkGreen}{rgb}{0,0.6,0}

\lstset{
	extendedchars=true,
	inputencoding=utf8,
	basicstyle=\small,
	keywordstyle=\color{red}\bfseries,
	identifierstyle=\color{darkGreen},
	commentstyle=\color{gray}\textit,
	stringstyle=\color{magenta},
	showstringspaces=false,
	tabsize=1,
	language=Haskell,
	numbers=left,         % Números na margem da esquerda
	numberstyle=\tiny,    % Números pequenos
	stepnumber=1,         % De 1 em 1
	numbersep=5pt         % Separados por 5pt
}

\lstset
{
	literate={ê}{{\^e}}1 {ç}{{\c{c}}}1 {ã}{{\~a}}1 {â}{{\^a}}1 {õ}{{\~o}}1 {í}{{\'i}}1 {é}{{\'e}}1 {á}{{\'a}}1 {ó}{{\'o}}1 {ú}{{\'u}}1
}

\title{Mini-Curso de Haskell - SEMAC 2013/1}
\author[M. Millani]{Marcelo Garlet Millani}
\institute[II/UFRGS]{Universidade Federal do Rio Grande do Sul \\ Instituto de Informática \\ Grupo PET Computação}

\begin{document}

	\begin{frame}
		\maketitle
	\end{frame}

\section{Sobre a Linguagem}

	\begin{frame}{}
		\tableofcontents[currentsection,currentsubsection]
	\end{frame}

	\begin{frame}{O que é Haskell}
		\begin{itemize}
		 \item Linguagem puramente funcional (não permite efeitos colaterais)
		 \begin{itemize}
		 \item ``Insanidade é continuar fazendo sempre a mesma coisa e esperar resultados diferentes''
		 
		 - Albert Einstein
		\end{itemize}
		
		 \item Usa o conceito de mônadas para isolar partes impuras do código
		 \item Usa avaliação preguiçosa (só calcula uma expressão se ela for necessária)
		\end{itemize}
		
	\end{frame}
	
	\begin{frame}{História}
	 
	 \begin{itemize}
	  \item Proposta na conferência FPCA '87 (Functional Programming Languages and Computer Architecture)
	  \item Havia diversas linguagens funcionais não-estritas, todas similares
	  \item Foi criada com o objetivo de propor uma linguagem funcional padrão
	  \item Nome em homenagem a Haskell Curry
	 \end{itemize}
	\end{frame}
	
	\begin{frame}{História}
	 A linguagem proposta deveria satisfazer o seguinte:
	 \begin{itemize}
	  \item Deveria ser adequada para o ensino, pesquisa e aplicações (incluindo sistemas grandes)
	  \item Deveria ser descrita completamente através da publicação de sintaxe e semânticas formais
	  \item Deveria estar disponível livremente, de forma que qualquer um pudesse implementar e distribuir a linguagem
	  \item Deveria ser baseada em ideias amplamente aceitas
	  \item Deveria reduzir diversidade desnecessária em linguagens funcionais
	 \end{itemize}
	\end{frame}
	
\section{Ferramentas}

	\begin{frame}{}
		\tableofcontents[currentsection,currentsubsection]
	\end{frame}

	\begin{frame}{Interpretadores e Compiladores}
	 
	 \begin{description}
	  \item [Hugs] Apenas um interpretador. É mais rápido que o interpretador do GHC;
			\begin{itemize}
			 \item \url{http://haskell.org/hugs}
			 \item Pacote hugs no linux
			\end{itemize}
	  \item [GHC] Compilador e interpretador. É o mais usado dos três;
			\begin{itemize}
			 \item \url{http://haskell.org/ghc}
			 \item Pacote ghc no linux
			\end{itemize}
	  \item [NHC] Apenas um compilador. Arquivos produzidos geralmente melhores que o do GHC.
			\begin{itemize}
			 \item \url{http://www.haskell.org/nhc98}
			\end{itemize}
	 \end{description}	 
	\end{frame}
	
	\begin{frame}{FREE HUGS!}
	 \begin{figure}
	  \includegraphics[height=0.8\textheight]{freeHugs.jpg}
	 \end{figure}

	\end{frame}

	
\section{Interpretador ghci}

		\begin{frame}{}
		 \tableofcontents[currentsection,currentsubsection]
		\end{frame}

	\begin{frame}[fragile]{Base}
	 	 \begin{itemize}
	 	  \item Declarações começam com \textit{let}
	 	 \end{itemize}
	 	 \begin{block}{Sintaxe}
	 	  let <identificador> {[}variáveis\dots{]} = <expressão>
	 	 \end{block}
	 	 \begin{block}{Exemplos}
	 	  \begin{lstlisting}
Prelude> let a = 12
Prelude> let soma x y = x+y
	 	  \end{lstlisting}
	 	 \end{block}
	\end{frame}
	
	\begin{frame}[fragile]{Controle de Fluxo}
		\begin{itemize}
		 \item Usaremos recursão em vez de laços
		 \item Possui construções do tipo if-then-else
		\end{itemize}
		
		\begin{block}{Exemplo}
		 \begin{lstlisting}
Prelude> if 2>1 then 3 else 9
		 \end{lstlisting}
		\end{block}
	\end{frame}
	
	\section{Operações}
	
	\begin{frame}{}
		 \tableofcontents[currentsection,currentsubsection]
	\end{frame}
	
	\begin{frame}{Operadores}
	 \begin{table}[h]
	  \centering
	  \begin{tabular}{c | c}
			\hline
			operador & operação\\
			\hline
			+ & soma \\
			- & subtração \\
			* & multiplicação\\
			\textasciicircum & potenciação\\
			/ & divisão (real)\\
			. & composição \\
			\$ & aplicação \\
			++ & concatenação de listas\\
			: & inserção em uma lista\\
			!! & acesso a uma posição da lista\\			
			\hline
	  \end{tabular}
	 \end{table}	 
	\end{frame}
	
	\begin{frame}[fragile]{Operações Inteiras}
	
	\begin{table}[h]
		\begin{tabular}{c | c}
			\hline
			função & operação\\
			\hline
			quot & divisão (truncada) \\
			div & divisão (floor) \\
			rem & resto \\
			mod & resto \\
			\hline
		\end{tabular}
	\end{table}
	
	\begin{block}{Propriedades}
		\begin{lstlisting}
Prelude> (quot x y)*y + (rem x y) == x  
Prelude> (div x y)*y + (mod x y) == x
		\end{lstlisting}
	 \end{block}

	\end{frame}
	
	\begin{frame}{Operadores Booleanos}
	 
	 \begin{table}[h]
	  \begin{tabular}{c | c}
	  \hline
	   operador & operação\\
	   \hline
	   \&\& & $\land$ \\
	   || & $\lor$ \\
	   not & $\lnot$ \\
	   \hline
	  \end{tabular}
	 \end{table}
	 
	 \begin{table}[h]
	  \begin{tabular}{c | c}
	  \hline
	   operador & operação\\
	   \hline
	   == & igual \\
	   /= & diferente \\
	   > & maior que \\
	   >= & maior ou igual \\
	   < & menor que \\
	   <= & menor ou igual \\
	   \hline
	  \end{tabular}
	 \end{table} 
	\end{frame}
	
	\section{Cálculo Lambda}
	
		\begin{frame}{}
		 \tableofcontents[currentsection,currentsubsection]
		\end{frame}

		\begin{frame}{Funções Lambda}
		 
		 \begin{itemize}
		  \item Funções sem nome
		  \item São avaliadas através de substituição de símbolos
		  \item Usa-se redução $\beta$ e $\alpha$ para fazer a avaliação
		 \end{itemize}
		 
		 \begin{block}{Exemplo}
		  \[(\lambda x. \lambda y. x+y)~ 5~ 4\]
		  \[(\lambda y. 5+y)~ 4\]
		  \[5+4\]		  
		 \end{block}		 
		\end{frame}
		
	\begin{frame}[fragile]{Funções Lambda}
	 
	 \begin{block}{Em Haskell}
	  \begin{lstlisting}	
Prelude> (\x -> x*x) 5
Prelude> (\x y ->x+y) 4 5
	  \end{lstlisting}	  
	 \end{block}
	\end{frame}

	\begin{frame}[fragile]{Currying}
		\begin{itemize}
			\item Avaliação parcial de funções
			\item Consequência direta do cálculo lambda
		\end{itemize}
		\begin{block}{Exemplo}
		 \begin{lstlisting}
Prelude> let a = \x y -> x+y
Prelude> let b = a 5
		 \end{lstlisting}
		\end{block}
	\end{frame}	
	
	\section{Tipos de Dados}
	
	\begin{frame}{}
		\tableofcontents[currentsection,currentsubsection]
	\end{frame}
	
		\begin{frame}[fragile]{Declaração}
			\begin{block}{Listas}
		\begin{lstlisting}
Prelude> [1,2,3,4]
Prelude> []
Prelude> [2,4..16]
Prelude> [1,3..] !! 5
		 \end{lstlisting}
		\end{block}
		
		\begin{block}{Tuplas}
		\begin{lstlisting}
Prelude> (1)
Prelude> (1,2,3)
Prelude> (1,"hello",[1,2])
Prelude> let tuple = (1 , (+) , 2)
		 \end{lstlisting}
		\end{block}
		
		\end{frame}
	
	\subsection{Listas}
	
	\begin{frame}[fragile]{Listas}
		\begin{block}{Map}
			\begin{lstlisting}
Prelude> map (*2) [1,3..11]
Prelude> map (\x -> x*x) [1,3..] !! 12
Prelude> filter (>3) [1,3,5,7]
			\end{lstlisting}
		\end{block}	 
	\end{frame}
	
	\begin{frame}[fragile]{Folding}
	 \begin{block}{foldr}
	  Associativo a direita
	  \begin{lstlisting}
Prelude> :t foldr
 foldr :: (a -> b -> b) -> b -> [a] -> b
Prelude> foldr (-) 0 [1,2,3,4]
Prelude> 1 - (2 - (3 - (4 - 0)))
	  \end{lstlisting}
	 \end{block}
	 
	 \begin{block}{foldl}
	  Associativo a esquerda
	  \begin{lstlisting}
Prelude> :t foldl
 foldl :: (a -> b -> a) -> a -> [b] -> a
Prelude> foldl (-) 0 [1,2,3,4]
Prelude> (((0 - 1) - 2) - 3) - 4
	  \end{lstlisting}
	 \end{block}
	\end{frame}
	
	\begin{frame}{Exercício}
	 \begin{block}{Exercício}
	  Usando fold e map, faça uma função para calcular a média aritmética e outra para o desvio padrão de uma lista de números. 
	 \end{block}
	 
	  \[ \overline{x} = \frac{\sum\limits_{i \in S}i}{\left|S\right|} \]
	  \[ \sigma^2 = \frac{\sum\limits_{i \in S}(i - \overline{x})^2}{\left|S\right|} \]	 

	\end{frame}
	
	\subsection{Determinando Tipos}
		\begin{frame}[fragile]{Determinando Tipos}
		 \begin{block}{Exemplo}
		 \begin{lstlisting}
Prelude> :t 5
 Num a => a
Prelude> :t "Hello"
 "Hello" :: [Char]
Prelude> :t (1,'a')
 Num t => (t, Char)
Prelude> :t (\x y -> x+y)
 Num a => a -> a -> a
		 \end{lstlisting}
		 \end{block}
		\end{frame}
	
		\begin{frame}[fragile]{Data}
		 
		 \begin{itemize}
		  \item Definição de estruturas
		 \end{itemize}
		 
		 \begin{block}{Definindo}
		  \begin{lstlisting}
Prelude> data Tree a = Leaf 
         | Node (Tree a) a (Tree a)
		  \end{lstlisting}
		 \end{block}		 
		 \begin{block}{Construindo}
		  \begin{lstlisting}
Prelude> let n2 = Node Leaf 10 Leaf
Prelude> let n1 = Node Leaf 6 Leaf
Prelude> let n0 = Node n1 7 n2
		  \end{lstlisting}
		 \end{block}
		\end{frame}
		
		\subsection{Pattern Matching}
		
		\begin{frame}[fragile]{Pattern Matching}
		\begin{itemize}
		 \item Comportamento com base na estrutura dos dados
		 \item Muito útil para tratar data types
		\end{itemize}	
		
		 \begin{block}{Exemplo}
		  \begin{lstlisting}
Prelude> let a = (\(x,y) -> x+y)
Prelude> let b = (\(x:y:_) -> x*y)
Prelude> let (p,k) = (\x y -> (x,y)) 3 4
Prelude> let Node esq valor dir = n0
		  \end{lstlisting}
		 \end{block}
		\end{frame}
		
		\subsection{Conversões}

			\begin{frame}[fragile]{Conversões de String}
			 \begin{block}{Funções}
			  \begin{lstlisting}
Prelude> show 32
Prelude> read "32"::Int
Prelude> read "32"::Float
			  \end{lstlisting}
			 \end{block}
			\end{frame}
			
			\begin{frame}[fragile]{Conversões de Números}
			 \begin{block}{Funções}
				\begin{lstlisting}
Prelude> :t fromIntegral
fromIntegral :: (Integral a, Num b) => a -> b
Prelude> :t toRational
Real a => a -> Rational
				\end{lstlisting}
			 \end{block}				
			\end{frame}
			
			\begin{frame}[fragile]{Aproximações}
			 \begin{block}{Funções}
			  \begin{lstlisting}
Prelude> :t ceiling
ceiling :: (Integral b, RealFrac a) => a -> b
Prelude> ceiling (-3.2)
Prelude> floor (-3.2)
Prelude> truncate (-3.2)
Prelude> round (-3.2)
			  \end{lstlisting}
			 \end{block}
			\end{frame}
			
			\begin{frame}{Aproximações}
			 \[ a \leq ceiling~ a\]
			 \[ floor~ a \leq a\]
			 \[ \left|truncate~ a\right| \leq \left|a\right|\]
			\end{frame}		
	
	\section{IO}
	
	\begin{frame}{}
		\tableofcontents[currentsection,currentsubsection]
	\end{frame}
	
		\begin{frame}[fragile]{Funções de Saída}
		 \begin{itemize}
		  \item Não são funções de fato
		  \item Permitem efeitos colaterais
		  \item São consideradas ``IO Actions''
		  \item Sempre retornam valores de um tipo IO
		  \item Usa-se o comando \textit{do} para executar uma sequência de ações
		 \end{itemize}
		 
		 \begin{block}{Print}
		  \begin{lstlisting}
print :: Show a => a -> IO ()
putStr :: String -> IO ()
putStrLn :: String -> IO ()
		  \end{lstlisting}
		 \end{block}
		\end{frame}
		
		\begin{frame}[fragile]{Funções de Saída}
		 \begin{block}{printing}
		  \begin{lstlisting}
Prelude> putStrLn "Nova linha"
Prelude> putStr "Mesma linha"
Prelude> print 42
		  \end{lstlisting}
		 \end{block}
		\end{frame}
		
		\begin{frame}[fragile]{Funções de Entrada}

		 \begin{block}{Definição}
		  \begin{lstlisting}
readFile :: FilePath -> IO String
getLine :: IO String
		  \end{lstlisting}
		 \end{block}
		 
		 \begin{block}{}
		  \begin{lstlisting}
Prelude> file <- readFile "arquivo"
Prelude> line <- getLine
		  \end{lstlisting}
		 \end{block}		 
		\end{frame}
		
		\begin{frame}[fragile]{mapM\_}
		\begin{itemize}
		 \item Similar ao map, mas para IO Actions
		\end{itemize}

		 \begin{block}{Exemplo}
		  \begin{lstlisting}
mapM_ (putStr . show) [1,4..10]
		  \end{lstlisting}
		 \end{block}

		\end{frame}


	\section{Arquivos}
	
	\begin{frame}{}
		\tableofcontents[currentsection,currentsubsection]
	\end{frame}
	
		\begin{frame}[fragile]{Arquivos}
		 \begin{itemize}
		  \item Código pode ser dividido em diversos arquivos
		  \item Permite a criação de módulos para facilitar reúso
		  \item Usa-se a extensão .hs
		 \end{itemize}
		 
		 \begin{block}{Exemplo}
		  \begin{lstlisting}
Module Main where
-- Comentario de uma linha
{-
	Comentario em bloco
-}
main = do
     putStrLn "Hello World"
		  \end{lstlisting}
		 \end{block}
		\end{frame}
		
		
		\begin{frame}[fragile]{Usando o GHC}
		 
		 \begin{block}{Compilação de arquivos específicos}		  
		 \begin{lstlisting}[language=Bash]
$ ghc arquivoA.hs arquivoB.hs
		 \end{lstlisting}
		 \end{block}
		 
		 \begin{block}{Compilação da main}
		  Nesse caso, o compilador busca as dependências (nome do arquivo deve ser o mesmo do módulo):
		  \begin{lstlisting}[language=Bash]
$ ghc --make arquivo.hs
			\end{lstlisting}
		 \end{block}
		 
		\end{frame}
		 
		 
		 \begin{frame}[fragile]{Algumas Opções do GHC}
		 
		 \begin{block}{Opções úteis}
		  \begin{itemize}
		   \item \verb+-o <ARQUIVO>+
				\begin{itemize}
				 \item Escolhe o arquivo de saída.
				\end{itemize}
		   \item \verb+-i<DIR1>[:<DIR2>...]+		   
				\begin{itemize}
				 \item Diretórios usados para buscar arquivos
				\end{itemize}
		   \item \verb+-odir <DIR>+
				\begin{itemize}
				 \item Diretório onde os arquivos objeto (.o) serão colocados
				\end{itemize}
			\item \verb+-hidir <DIR>+
				\begin{itemize}
				 \item Diretório onde os arquivos interface (.hi) serão colocados
				\end{itemize}
		  \end{itemize}
		 \end{block}
		 
		\end{frame}		
		
		\begin{frame}[fragile]{Definição de Funções}
		 \begin{block}{Formato}
		  \begin{lstlisting}
<identificador> [argumentos...] = <expressão>
		  \end{lstlisting}
		 \end{block}
		 
		 \begin{block}{Exemplo}
		  \begin{lstlisting}
fatorial 0 = 1
fatorial 1 = 1
fatorial n = n * fatorial (n-1)
		  \end{lstlisting}
		 \end{block}
		\end{frame}
		
		\begin{frame}[fragile]{Function Guards}
		 \begin{itemize}
		  \item Sintaxe alternativa para separação em casos
		  \item \textit{otherwise} se comporta como um \textit{wildcard}
		 \end{itemize}
		 \begin{block}{Exemplo}
		 \begin{lstlisting}
fatorial n
  | n==0 = 1
  | n==1 = 1
  | otherwise = n * fatorial (n-1)
		\end{lstlisting}
		 \end{block}
		\end{frame}

		
		\begin{frame}[fragile]{Let e Where}
		 \begin{itemize}
		  \item Permite a definição de variáveis auxiliares
		  \item Restringem o escopo
		  \item Múltiplas definições devem estar alinhadas
		 \end{itemize}
		 \begin{block}{Let-in}
		  \begin{lstlisting}
variance l =
 let avg = average l
     dif = map (\x -> (x - avg)^2) l
 in average dif
		  \end{lstlisting}	  
		 \end{block}
		\end{frame}
		
		\begin{frame}[fragile]{Where}
		 \begin{block}{Where}
		  \begin{lstlisting}
variance l =
 average dif
 where avg = average l
       dif = map (\x -> (x - avg)^2) l
		  \end{lstlisting}
		 \end{block}		 
		\end{frame}
		
		\begin{frame}{Exercício}
		 \begin{block}{Exercício}
		  Faça um programa para produzir uma lista de vértices a partir de uma forma geométrica básica e um ponto (para círculos e retângulos).
		 \end{block}
		 \begin{description}
		  \item [Círculo] Raio
		  \item [Retângulo] Altura e largura
		 \end{description}
		 Para discretizar um círculo de raio $r$ centrado em $(0,0)$:
		 \[
		  P_k = ( r \cdot sen(\frac{2 k \pi}{n}) , r \cdot cos(\frac{2 k \pi}{n} ) ) , \forall k \in [1,n]
		 \]


		\end{frame}
		
		\subsection{Módulos}
			\begin{frame}[fragile]{Módulos}
			 \begin{itemize}
			  \item Servem para organizar o código
			  \item Permitem reúso de código
			  \item Fáceis de se usar
			 \end{itemize}
			\end{frame}
			
			\begin{frame}[fragile]{Definição}
			 \begin{block}{Exemplo}
			  \lstinputlisting{exemplos/modules/Arvore.hs}
			 \end{block}
			\end{frame}
			
			\begin{frame}[fragile]{Import}
			 \begin{block}{Variações}
			  \begin{description}
			   \item [import Module] traz todas as funções de Module
					\begin{itemize}
					 \item x, y, z, Module.x, Module.y, Module.z
					\end{itemize}
				 \item [import Module (x,y)] traz somente x e y
				 \item [import Module as Foo] trata Module como Foo
					\begin{itemize}
					 \item x, y, z, Foo.x, Foo.y, Foo.z
					\end{itemize}

			   \item [import qualified Module] requer o nome do módulo prefixado
					\begin{itemize}
					 \item Module.x, Module.y, Module.z
					\end{itemize}
			  \end{description}			  
			 \end{block}
			\end{frame}
			
			\begin{frame}[fragile]{Import}
			 \begin{block}{Exemplo}
			  \lstinputlisting{exemplos/src/arvore.hs}
			 \end{block}
			\end{frame}
			
			\begin{frame}{Exercício}
			 \begin{block}{Exercício}
			  Transforme o último exercício em um módulo e use-o em outro arquivo.
			 \end{block}
			\end{frame}

			
\section{Referências}
\begin{frame}{Referências}
	\begin{thebibliography}{9}
	\bibitem{YAHT} Hal Daumé III, \textit{Yet Another Haskell Tutorial}
	\bibitem{THSE} Hudak, Paul. \textit{The Haskell School of Expression}
	\end{thebibliography}
\end{frame}
	
	
\end{document}
